% Part: first-order-logic
% Chapter: tableaux
% Section: quantifier-rules

\documentclass[../../../include/open-logic-section]{subfiles}

\begin{document}

\olfileid{fol}{tab}{qrl}

\olsection{قواعد سورها}

\subsection{قواعد $\lforall$}

\begin{defish}
\AxiomC{\sFmla{\True}{\lforall[x][!A(x)]}}
\RightLabel{\TRule{\True}{\forall}}
\UnaryInfC{\sFmla{\True}{!A(t)}}
\DisplayProof
\hfill
\AxiomC{\sFmla{\False}{\lforall[x][!A(x)]}}
\RightLabel{\TRule{\False}{\lforall}}
\UnaryInfC{\sFmla{\False}{!A(a)}}
\DisplayProof
\end{defish}

در \TRule{\True}{\lforall}، $t$ یک ترم بسته است (یعنی ترمی بدون
متغیر). در \TRule{\False}{\lforall}، $a$~!!a{constant} است که نباید در
هیچ‌جای شاخهٔ بالای قاعدهٔ \TRule{\False}{\lforall} رخ دهد. $a$ را
\emph{متغیر ویژهٔ} استنتاج \TRule{\False}{\forall} می‌نامیم.\footnote{هرچند
$a$ در قاعدهٔ بالا !!a{constant} است، اصطلاح «متغیر ویژه» را به کار
می‌بریم. این کاربرد دلایل تاریخی دارد.}

\subsection{قواعد $\lexists$}

\begin{defish}
\AxiomC{\sFmla{\True}{\lexists[x][!A(x)]}}
\RightLabel{\TRule{\True}{\lexists}}
\UnaryInfC{\sFmla{\True}{!A(a)}}
\DisplayProof
\hfill
\AxiomC{\sFmla{\False}{\lexists[x][!A(x)]}}
\RightLabel{\TRule{\False}{\lexists}}
\UnaryInfC{\sFmla{\False}{!A(t)}}
\DisplayProof
\end{defish}

باز هم $t$ یک ترم بسته است و $a$~!!a{constant} است که در شاخهٔ بالای
قاعدهٔ~\TRule{\True}{\lexists} رخ نمی‌دهد. $a$ را \emph{متغیر ویژهٔ}
استنتاج \TRule{\True}{\lexists} می‌نامیم.

شرطِ رخ‌ندادن متغیر ویژه در شاخهٔ بالای استنتاج
\TRule{\False}{\lforall} یا \TRule{\True}{\lexists}، \emph{شرط متغیر ویژه}
نامیده می‌شود.

\begin{explain}
قرارداد را به یاد آورید که هرگاه $!A$~!!a{formula} با !!{variable} آزادِ
~$x$ باشد، این امر را با نوشتن~$!A(x)$ نشان می‌دهیم. در همان بافت،
$!A(t)$ اختصاری برای~$\Subst{!A}{t}{x}$ است. بنابراین می‌توانستیم قاعدهٔ
\TRule{\False}{\lexists} را به صورت زیر نیز بنویسیم:
\begin{prooftree}
  \AxiomC{\sFmla{\False}{\lexists[x][!A]}}
  \RightLabel{\TRule{\False}{\lexists}}
  \UnaryInfC{\sFmla{\False}{\Subst{!A}{t}{x}}}
\end{prooftree}
توجه کنید که $t$ ممکن است از پیش در~$!A$ رخ دهد؛ برای نمونه، ممکن است
$!A$ همان~$\Atom{\Obj P}{t,x}$ باشد. پس استنتاج
$\sFmla{\False}{\Atom{\Obj P}{t,t}}$ از
~$ \sFmla{\False}{\lexists[x][\Atom{\Obj P}{t,x}]}$ کاربردی درست از
~\TRule{\False}{\lexists} است. با این حال، شروط متغیر ویژه در
\TRule{\False}{\lforall} و~\TRule{\True}{\lexists} ایجاب می‌کنند که
!!{constant}~$a$ در~$!A$ رخ ندهد. بنابراین نمی‌توان
$\sFmla{\False}{\Atom{\Obj P}{a,a}}$ را از
$\sFmla{\False}{\lforall[x][\Atom{\Obj P}{a,x}]}$ با استفاده
از~$\TRule{\False}{\lforall}$ به‌درستی استنتاج کرد.
\end{explain}

\begin{explain}
در \TRule{\True}{\lforall} و \TRule{\False}{\lexists} هیچ محدودیتی بر
ترم~$t$ وجود ندارد. از سوی دیگر، در قواعد \TRule{\True}{\lexists} و
\TRule{\False}{\lforall}، شرط متغیر ویژه ایجاب می‌کند که !!{constant}~$a$
در هیچ‌جای شاخه‌های بالای استنتاج مربوط رخ ندهد. این شرط برای تضمین صحت
دستگاه ضروری است. بدون آن، عبارت زیر !!{tableau}ی بسته برای
$\lexists[x][\formula{A}(x)] \lif \lforall[x][\formula{A}(x)]$ می‌بود:
\begin{center}
\begin{tableau}{}
  [\sFmla{\False}{\lexists[x][\formula{A}(x)] \lif \lforall[x][\formula{A}(x)]}, just=\TAss
    [\sFmla{\True}{\lexists[x][\formula{A}(x)]},
      just={\TRule{\False}{\lif}[1]}
      [\sFmla{\False}{\lforall[x][\formula{A}(x)]},
        just={\TRule{\False}{\lif}[1]}
        [\sFmla{\True}{\formula{A}(a)},
          just={\TRule{\True}{\lexists}[2]}
          [\sFmla{\False}{\formula{A}(a)},
            just={\TRule{\False}{\lforall}[3]}, close]
        ]
      ]
    ]
  ]
\end{tableau}
\end{center}
با این همه، $\lexists[x][\formula{A}(x)] \lif
\lforall[x][\formula{A}(x)]$ معتبر نیست.
\end{explain}

\end{document}
